%% theory_draft.tex
%% Standalone theory section for review — not yet integrated into main draft.
%% Flag to Neil before merging: check notation consistency with rest of paper.

\section{Theory}
\label{sec:theory}

We develop a model of investor belief updating around public audit signals to
formalize how political polarization affects market outcomes. The model draws on
the heterogeneous-priors framework of \citet{kandel1995} and the polarization
measure of \citet{esteban1994}.

\subsection{Setting}

Consider a firm whose true quality $\theta \in \{0, 1\}$ is unknown to
investors. The firm discloses an auditor-related signal $s \in \{0, 1\}$ (e.g.,
an auditor change announcement). Signals are informative: $P(s=1 \mid
\theta=1) = p_H > P(s=1 \mid \theta=0) = p_L$, with $p_H > p_L$.

A continuum of investors $i \in [0,1]$ observe $s$ and form posterior beliefs.
Investors share the same likelihood function $\{p_H, p_L\}$ but differ in their
prior beliefs $\pi_i \equiv P_i(\theta = 1) \in (0,1)$.

\textbf{Source of prior heterogeneity.} We model investors as belonging to
ideological groups indexed by political affiliation $g \in \{D, R\}$, with
population shares $s_D$ and $s_R = 1 - s_D$. Group-level priors reflect
differential trust in auditors and regulatory institutions:
$\pi_D \neq \pi_R$. Polarization increases the \emph{dispersion} of priors
across groups without changing the \emph{average} prior.

\subsection{Belief Updating}

By Bayes' rule, investor $i$'s posterior belief after observing $s=1$ is:

\begin{equation}
    \label{eq:bayes}
    P_i(\theta = 1 \mid s = 1)
    = \frac{p_H \, \pi_i}{p_H \, \pi_i + p_L (1 - \pi_i)}
    \equiv \phi(\pi_i).
\end{equation}

The function $\phi(\cdot)$ is strictly increasing and concave in $\pi_i$.
Investors with higher priors update to higher posteriors after a positive
signal, but the updating function compresses the distribution of posteriors
relative to the distribution of priors.

\textbf{Effect of polarization on posterior dispersion.}
Let $\Pi \equiv \{\pi_i\}_{i \in [0,1]}$ denote the cross-sectional
distribution of priors. Define posterior dispersion as
$\sigma^2_{\text{post}} \equiv \text{Var}_i[\phi(\pi_i)]$.
Because $\phi$ is monotone, greater dispersion in $\Pi$ generates greater
dispersion in posteriors. Formally:

\begin{proposition}[Polarization and belief dispersion]
    \label{prop:dispersion}
    Let $\Pi'$ be a mean-preserving spread of $\Pi$ (i.e., $\Pi'$ is more
    dispersed in the sense of second-order stochastic dominance).
    Then $\text{Var}_i[\phi(\pi_i) \mid \Pi'] \geq \text{Var}_i[\phi(\pi_i)
    \mid \Pi]$.
\end{proposition}

\begin{proof}
    Since $\phi$ is monotone increasing, the composition $\phi \circ \Pi$ is
    a monotone transformation. For any monotone function, a mean-preserving
    spread in the argument weakly increases the variance of the image.
    $\square$
\end{proof}

\subsection{Equilibrium Prices and Trading}

Following \citet{harris1993}, investors trade to maximize expected utility.
In equilibrium, the asset price aggregates beliefs:

\begin{equation}
    \label{eq:price}
    P^* = \int_0^1 w_i \, E_i[\theta \mid s] \, di
        = \int_0^1 w_i \, \phi(\pi_i) \, di,
\end{equation}

where $w_i$ are equilibrium portfolio weights that depend on risk tolerance
and endowments. For simplicity, assume equal weights ($w_i = 1$ for all $i$).

Trading volume arises from disagreement: investor $i$ trades if and only if
$\phi(\pi_i) \neq P^*$. Total trading volume is proportional to the cross-
sectional dispersion of posterior beliefs \citep{kandel1995}:

\begin{equation}
    \label{eq:volume}
    V \propto \text{Var}_i[\phi(\pi_i)].
\end{equation}

\subsection{Role of Signal Ambiguity}

When the signal is more ambiguous — that is, when $p_H - p_L$ is smaller —
the likelihood function provides less information, and the posterior
$\phi(\pi_i)$ is more sensitive to the prior $\pi_i$. Formally,
$\partial^2 \phi / \partial \pi_i \, \partial (p_H - p_L) < 0$: a smaller
signal precision amplifies the role of the prior in determining the posterior.
Consequently, polarization effects on disagreement are stronger when signals
are less informative.

\subsection{Testable Predictions}

The model yields three empirical predictions, which map directly to our
empirical specification:

\begin{itemize}
    \item[\textbf{P1.}] \textbf{Absolute returns.}
        Greater prior dispersion (higher polarization) generates greater
        dispersion in posteriors. When investors trade at different prices,
        the equilibrium price reflects a weighted average of heterogeneous
        valuations, amplifying both the magnitude and cross-sectional
        variance of returns. Therefore, polarization increases $|CAR|$.

    \item[\textbf{P2.}] \textbf{Trading volume.}
        By equation~(\ref{eq:volume}), trading volume is proportional to
        posterior belief dispersion. Since polarization increases dispersion
        (Proposition~\ref{prop:dispersion}), polarization increases abnormal
        trading volume around audit signal events.

    \item[\textbf{P3.}] \textbf{Signal ambiguity.}
        The effects in P1 and P2 are stronger for more ambiguous audit signals
        (smaller $p_H - p_L$). In the empirical analysis, we proxy for
        ambiguity using auditor change type: dismissals without stated reasons
        and switches between auditors of similar size are more ambiguous than
        Big-4-to-non-Big-4 switches with disclosed disagreements.
\end{itemize}

\textbf{Discussion.} The model predicts that polarization affects the
\emph{processing} of identical information, not the information itself.
This distinguishes our mechanism from alternative explanations based on
differential information access or manipulation of audit quality. In our
setting, the signal $s$ is identical across all investors; polarization
operates solely through the prior $\pi_i$.

This distinction has regulatory implications. Standard disclosure-based
interventions (e.g., expanded auditor reporting requirements) cannot
eliminate polarization-driven disagreement, because the source of
heterogeneity is in investors' priors rather than in the information set.
